\section{Od sfery do pełnej kuli}
\subsection{Przedłużenie kawałków wzdłuż promieni}
Niech $B^3_*=B^3\setminus\{0\}$. Dla $H\subseteq S^2$ definiujemy jego stożek promieniowy w kuli bez środka:
\[
 \widehat H=\{rx:x\in H,\ 0<r\le1\}.
\]
Każdy $y\in B^3_*$ ma dokładnie jedno przedstawienie $y=rx$ z $0<r\le1$, $x\in S^2$: mianowicie $r=\|y\|$, $x=y/\|y\|$. Dlatego partycja sfery daje partycję $B^3_*$ przez odpowiednie stożki.
\begin{lemat}[Obroty komutują z przedłużeniem promieniowym]\label{lem:stozki}
Dla $Q\in\mathrm{SO}(3)$ i $H\subseteq S^2$ mamy $Q\widehat H=\widehat{QH}$.
\end{lemat}
\begin{proof}
Liniowość macierzy daje $Q(rx)=rQx$, a ortogonalność daje $\|Qx\|=1$. To dowodzi jednej inkluzji. Jeśli $ry$ należy do $\widehat{QH}$, to $y=Qx$ dla pewnego $x\in H$, więc $ry=Q(rx)\in Q\widehat H$. Mamy także drugą inkluzję.
\end{proof}
\begin{wniosek}\label{wn:kula-bez-srodka}
$B^3_*\equiv2B^3_*$ za pomocą co najwyżej $16$ kawałków i obrotów.
\end{wniosek}
\begin{proof}
Bierzemy stożki nad częściami $H_{ijk}$ z \eqref{eq:7-2} i stosujemy te same obroty $q_{ijk}$. Jednoznaczność kierunku i promienia zapewnia rozłączność części wyjściowych, a lemat~\ref{lem:stozki} przenosi rozłączność i pokrycie obrazów z każdej sfery na całą kulę bez środka. Nie zmieniamy promienia żadnego punktu podczas tego etapu.
\end{proof}
\begin{uwaga}[Dlaczego nie dopisać po prostu środka?]
Wszystkie obroty wokół $0$ ustalają $0$. Z jednego takiego punktu nie otrzymamy dwóch punktów oznaczonych różnymi etykietami przez dopisanie go do dwóch części --- przestałyby być rozłączne. Środek wymaga osobnego argumentu i ruchu wokół osi nieprzechodzącej przez $0$.
\end{uwaga}

\subsection{Wchłonięcie środka za pomocą obrotu}
W płaszczyźnie $z=0$ weźmy okrąg
\[
 \Gamma=\{(1/4,0,0)+(1/4)(\cos t,\sin t,0):t\in\R\}.
\]
Ma środek $c=(1/4,0,0)$ i promień $1/4$. Punkt $0$ należy do $\Gamma$ (dla $t=\pi$). Każdy punkt okręgu ma normę nie większą niż $\|c\|+1/4=1/2$, więc cały okrąg leży wewnątrz $B^3$.

Niech $U$ będzie obrotem wokół osi $z$ o kąt $\alpha=2\pi\sqrt2$. Określmy ruch sztywny
\begin{equation}\label{eq:8-1}
 T_0(x)=c+U(x-c)=Ux+(c-Uc).
\end{equation}
Jest to obrót wokół prostej $c+\R(0,0,1)$; zachowuje okrąg $\Gamma$.
\begin{lemat}\label{lem:wchloniecie-srodka}
Punkty $0,T_0(0),T_0^2(0),\ldots$ są parami różne. W konsekwencji $B^3\equiv B^3_*$ za pomocą dwóch części i ruchów sztywnych zachowujących orientację.
\end{lemat}
\begin{proof}
Gdyby $T_0^m(0)=T_0^n(0)$ dla $m>n$, to $T_0^{m-n}(0)=0$. Punkt $0$ leży na okręgu o dodatnim promieniu, więc wraca na swoje miejsce dokładnie wtedy, gdy kąt obrotu jest całkowitą wielokrotnością $2\pi$. Dostajemy $(m-n)\alpha=2\pi k$, a więc $(m-n)\sqrt2=k$ dla pewnego $k\in\Z$. To sprzeczność z niewymiernością $\sqrt2$.

Dla kompletności: gdyby $\sqrt2=p/q$ w postaci nieskracalnej, to $p^2=2q^2$, więc $p$ byłoby parzyste. Zapis $p=2r$ dawałby $q^2=2r^2$, więc także $q$ byłoby parzyste, wbrew nieskracalności.

Niech $M=\{T_0^n(0):n\ge0\}$. Wszystkie te punkty leżą na $\Gamma\subset B^3$. Stosujemy lemat~\ref{lem:lancuch} dla $D=\{0\}$, $E=B^3$, $T=T_0$. Częściami są $M$ i $B^3\setminus M$, a ruchami $T_0$ i identyczność. Ich obrazy to $M\setminus\{0\}$ oraz $B^3\setminus M$, czyli partycja $B^3_*$.
\end{proof}
\begin{uwaga}
Nie twierdzimy, że $T_0(B^3)=B^3$. Nie jest to potrzebne. Obracamy wyłącznie kawałek $M$, który leży na zachowywanym okręgu. Reszta kuli nie jest poruszana. Dowolny ruch nie musi zachowywać całej kuli, tylko poprawnie przenosić przypisany mu kawałek.
\end{uwaga}

\subsection{Złożenie wszystkich etapów}
\begin{proof}[Dowód twierdzenia~\ref{tw:banach-tarski}]
Lemat~\ref{lem:wchloniecie-srodka} i wniosek~\ref{wn:kula-bez-srodka} dają $B^3\equiv B^3_*\equiv2B^3_*\equiv2B^3$. Przechodniość została udowodniona przez jawne rozdrobnienie części w lemacie~\ref{lem:skladanie}. Nie używamy więc żadnego nieskończonego składania ruchów.

Podajmy również bezpośredni wzór na końcowe części. Zanumerujmy niepuste części rozkładu sfery jako $H_v$, ich obroty jako $q_v$, a etykiety jako $\ell(v)\in\{1,2\}$, gdzie $v$ przebiega zbiór najwyżej $16$ indeksów. Niech $\widehat H_v$ oznacza ich stożki promieniowe. Przyjmijmy $J_0=M$, $J_1=B^3\setminus M$, $t_0=T_0$, $t_1=\id$. Bijekcja $k:B^3\to B^3_*$ działa na $J_s$ jako $t_s$. Zbiory $t_0J_0$ i $t_1J_1$ tworzą partycję $B^3_*$.

Dla $s,t\in\{0,1\}$ i indeksu $v$ definiujemy
\begin{equation}\label{eq:8-2}
 P_{svt}=\{x\in J_s:t_sx\in\widehat H_v,\quad q_vt_sx\in t_tJ_t\}.
\end{equation}
Na tym zbiorze złożenie bijekcji $k$, bijekcji wyznaczonej przez paradoksalny rozkład $B^3_*$ i odwrotnej bijekcji $k^{-1}$ ma postać jednego ruchu
\begin{equation}\label{eq:8-3}
 x\longmapsto t_t^{-1}q_vt_sx,
\end{equation}
z etykietą $\ell(v)$. Dla każdego $x\in B^3$ istnieje jedyne $s$ z $x\in J_s$, jedyne $v$ z $t_sx\in\widehat H_v$ i jedyne $t$ określające, w której części $t_tJ_t$ leży $q_vt_sx$. To dowodzi, że \eqref{eq:8-2} jest partycją całej $B^3$.

Ponieważ każdy etap jest bijekcją między odpowiednimi zbiorami etykietowanymi, obrazy pod ruchami \eqref{eq:8-3} tworzą partycję dwóch formalnych kul. Liczba części jest nie większa niż $2\cdot16\cdot2=64$. Nie ma dodatkowego czynnika $2$ za etykiety: każdy indeks $v$ ma już jedną, określoną etykietę docelową.

Teraz umieszczamy te kopie w przestrzeni. Dla etykiety $1$ dodajemy $v_1=(-3,0,0)$, a dla etykiety $2$ dodajemy $v_2=(3,0,0)$. Ostateczny ruch dla części $P_{svt}$ to $f_{svt}(x)=t_t^{-1}q_vt_sx+v_{\ell(v)}$. Każde $t_s$ jest ruchem $x\mapsto Qx+v$ z $Q\in\mathrm{SO}(3)$, a każde $q_v$ obrotem. Odwrotności i złożenia takich ruchów znów mają tę postać: jeśli $f(x)=Qx+v$, $g(x)=Rx+w$, to $g(f(x))=RQx+(Rv+w)$ oraz $f^{-1}(x)=Q^{-1}x-Q^{-1}v$.

Zatem wszystkie $f_{svt}$ są dozwolonymi ruchami sztywnymi. Ich obrazy wypełniają dokładnie dwie wskazane kule i są parami rozłączne: wewnątrz każdej kuli wynika to z bijekcji, a między kulami z ich rozłącznego położenia. To kończy dowód.
\end{proof}

\subsection{Kula o innym środku i promieniu}
\begin{wniosek}\label{wn:dowolna-kula}
Każdą kulę domkniętą o promieniu $r>0$ można podzielić na skończenie wiele części, a następnie za pomocą ruchów sztywnych zachowujących orientację złożyć z nich dwie rozłączne kule tego samego promienia.
\end{wniosek}
\begin{proof}
Dla kuli o środku $c$ użyjmy pomocniczego podobieństwa $S(x)=c+rx$ do zdefiniowania części $S(P_i)$ z dowodu dla kuli jednostkowej. Jeśli ruch jednostkowy miał postać $f_i(x)=Q_ix+v_i$, to $Sf_iS^{-1}(y)=Q_iy+(c-Q_ic+rv_i)$. To nadal izometria: liniowa część wynosi $Q_i$, nie $rQ_i$. Podobieństwo służy tylko przeniesieniu opisu konstrukcji; rzeczywiste kawałki nowej kuli poruszamy powyższymi izometriami, bez skalowania.
\end{proof}

